\section{模型评价与推广}
本文四个问题始终围绕同一微网物理结构逐步扩展，问题一建立基础电能流和储能状态关系，问题二将全天计划与实际执行分开，问题三增加日内新预报到达后的再决策，问题四进一步把实时电价作为随机量。这样处理的优点在于前一问的变量和约束可以直接被后一问继承，模型之间不存在明显割裂，也便于解释每一项新增变量究竟对应哪一条题意。

基础模型将外网购电区分为直接供负荷和用于储能充电，将富余光伏区分为储能充电和弃光，并在问题二以后再用 \(W\) 单独描述已购未用电，因此能够较清楚地追踪外网、光伏和储能之间的实际流向。问题二以后只把真正需要提前提交的计划购电作为事前变量，而把充放电、弃光和紧急购电留到情景实现后决定，较好地符合真实信息时序。利用历史完整日预测误差构造随机情景，可以避免把单一预测曲线当成未来真值；7 日滚动视野又使当天日末储电量具有未来价值，减轻单日模型在时域末端过度放电的问题。

问题三利用 0:00、6:00、12:00、18:00 新光伏预报重新调整尚未执行计划，并始终使用真实执行后的储电量连接不同阶段，使“新信息—再决策—实际执行”形成连续链条。问题四保留负价格情景而不是简单截断，通过计划购电物理上界和调增、调减互斥约束消除非物理套利，这一处理能够在不改变原结算规则的前提下保证模型有界。

模型仍存在一定局限。当前目标函数主要考虑购电费用，没有显式加入储能循环衰减和设备寿命损耗，因此频繁充放电的长期成本可能被低估。正式随机模型的情景数受全年滚动计算量限制，问题二结果已经表明紧急购电量对情景覆盖较敏感；问题三中日内只更新光伏预报，负荷中心预测仍沿用 0:00 版本；问题四中调度效果又会受到价格峰谷时刻预测能力限制。与此同时，本文按照题面没有设置余电上网和需求响应机制，所以适用范围主要是“负荷—光伏—储能—外网购电”这一类微网调度。

后续可从几个方向继续改进。首先，可在目标函数中增加与充放电通量和循环深度相关的储能衰减成本，使短期电费与长期设备寿命同时进入优化。其次，可以增加随机情景数量，并配合聚类或情景削减方法控制计算规模。若计算资源允许，还可以把当前阶段滚动随机规划进一步扩展为严格的多阶段随机规划，使每个阶段的决策只依赖当时已经观察到的信息。问题三和问题四中还可以利用已发生负荷对当天剩余负荷进行日内修正，并进一步提高价格峰谷时刻预测能力。

从推广角度看，本文基础物理层由负荷、分布式电源、储能和外部电网组成，只要重新给定容量、功率、效率、电价和预测误差数据，同一思路还可以用于园区微网、商业建筑群、校园能源系统以及含风电等其他随机可再生能源的储能调度。整个框架可以理解为“物理能量关系 + 预测 + 随机情景 + 滚动决策 + 市场结算”几个模块的组合，具体应用场景变化时可以替换其中某一模块，而不必重新建立全部模型。
